% !TEX root = 概率论期中复习.tex
\clearpage\section{第二周：概率空间、概率界与事件极限}
\subsection{教材第19页第1题}\label{ex:w02-t1}
\paragraph{题目}
假设某个随机试验包含可数无穷多个不同的结果，且每个结果出现的概率均严格大于零。解释为何不会所有结果发生的概率都相同。
\basis{原笔记第1.2节概率的可列可加性。}
\paragraph{思路}若各个结果具有同一个正概率，有限多个互斥结果的总概率也将超过1。
\begin{proof}[解答]
把这些结果列为$\omega_1,\omega_2,\ldots$，单点事件两两互斥。若所有单点的概率同为$p>0$，则对任意正整数$N$，
\[P(\{\omega_1,\ldots,\omega_N\})=Np\leq1.\]
取$N>1/p$即矛盾。因此正概率不能全部相同。若试图令共同值为0，可列可加性又给出$P(\Omega)=\sum_{i=1}^\infty0=0$，也不可能。
\end{proof}
\paragraph{答案}不存在可数无穷等可能且每个结果有正概率的模型。

\subsection{教材第19页第2题(1)}\label{ex:w02-t2}
\paragraph{题目}
设$(\Omega,\mathscr F)$为可测空间。给定非空集$E\subseteq\Omega$，记$\mathscr E=\{E\cap A:A\in\mathscr F\}$。证明$\mathscr E$是$E$上的$\sigma$代数。（按本周要求不做第(2)小题。）
\basis{原笔记第1.2节“可测空间”中$\sigma$代数的定义。}
\paragraph{思路}以$E$为新的全集，分别验证包含全集、相对补封闭和可列并封闭。
\begin{proof}
因为$\Omega\in\mathscr F$，有$E=E\cap\Omega\in\mathscr E$。

若$B\in\mathscr E$，则存在$A\in\mathscr F$使$B=E\cap A$。在新全集$E$中，$B$的补集为
\[E\setminus B=E\setminus(E\cap A)=E\cap A^c\in\mathscr E,\]
这里用了$A^c\in\mathscr F$。

若$B_i\in\mathscr E$，分别写成$B_i=E\cap A_i$，其中$A_i\in\mathscr F$，则
\[\bigcup_{i=1}^{\infty}B_i=E\cap\left(\bigcup_{i=1}^{\infty}A_i\right)\in\mathscr E,\]
因为$\bigcup_iA_i\in\mathscr F$。三个条件均满足，故$\mathscr E$是$E$上的$\sigma$代数。
\end{proof}
\paragraph{答案}$\mathscr E$是$E$上的$\sigma$代数，补运算必须相对于$E$进行。

\subsection{教材第19页第4题}\label{ex:w02-t4}
\paragraph{题目}
设$(\Omega,\mathscr F,P)$为概率空间。证明对$A,B\in\mathscr F$有
\[P(A\cap B)\geq P(A)+P(B)-1.\]
\basis{原笔记第1.2节命题“一般加法公式”。}
\paragraph{思路}用并事件的概率不超过1。
\begin{proof}
由加法公式，
\[P(A\cap B)=P(A)+P(B)-P(A\cup B).\]
因为$A\cup B\subseteq\Omega$，单调性给出$P(A\cup B)\leq1$。代入即得
\[P(A\cap B)\geq P(A)+P(B)-1.\]
\end{proof}
\paragraph{答案}所给下界成立；结合非负性还可写成$P(AB)\geq\max\{0,P(A)+P(B)-1\}$。

\subsection{教材第19页第5题}\label{ex:w02-t5}
\paragraph{题目}
设事件$A,B$满足$P(A)=3/4$和$P(B)=1/3$。证明$1/12\leq P(A\cap B)\leq1/3$。
\basis{原笔记第1.2节单调性、加法公式；\hyperref[ex:w02-t4]{上一题}。}
\paragraph{思路}下界用$P(A)+P(B)-1$，上界用$A\cap B\subseteq B$。
\begin{proof}
由上一题，
\[P(A\cap B)\geq\frac34+\frac13-1=\frac1{12}.\]
另一方面，$A\cap B\subseteq B$，所以$P(A\cap B)\leq P(B)=1/3$，合并即得结论。
\end{proof}
\paragraph{答案}$P(A\cap B)\in[1/12,1/3]$。

\subsection{作业第2题：偏离独立乘积的界}\label{ex:w02-h2}
\paragraph{题目}
求证：对任意事件$A,B\in\mathscr F$，有
\[|P(AB)-P(A)P(B)|\leq\frac14.\]
提示：不妨设$P(A)\leq P(B)$。注意不等式$ab\leq((a+b)/2)^2$。
\basis{原笔记第1.2节概率的单调性、加法公式。}
\paragraph{思路}按差值的正负分别估计，并按提示设$a=P(A)\leq b=P(B)$。
\begin{proof}
记$q=P(AB)$。先估计正向差值。由$q\leq a$与$a\leq b$，有
\begin{align*}
q-ab&\leq a(1-b)\leq a(1-a)\\
&=\frac14-\left(a-\frac12\right)^2\leq\frac14.
\end{align*}
再估计反向差值。若$a+b\leq1$，由$q\geq0$及提示中的不等式，
\[ab-q\leq ab\leq\left(\frac{a+b}{2}\right)^2\leq\frac14.\]
若$a+b>1$，由$q\geq a+b-1$，得到
\begin{align*}
ab-q&\leq ab-a-b+1=(1-a)(1-b)\\
&\leq\left(\frac{2-a-b}{2}\right)^2<\frac14.
\end{align*}
两种情况穷尽全部可能，故$-1/4\leq q-ab\leq1/4$，结论成立。
\end{proof}
\paragraph{答案}$1/4$为可达到的界，例如$A=B$且$P(A)=1/2$时取等号。

\subsection{作业第3题：开区间的可列表示}\label{ex:w02-h3}
\paragraph{题目}
求证：对任意实数$a<b$，有
\[(a,b)=\bigcup_{k=1}^{\infty}(a,b-1/k].\]
\basis{原笔记第1.2节常用博雷尔集合与阿基米德性质。}
\paragraph{思路}先证每个右侧区间包含于左侧，再对左侧每个点选取足够大的$k$。
\begin{proof}
若$x\in(a,b-1/k]$，则$a<x\leq b-1/k<b$，所以右侧并集包含于$(a,b)$。这里当$b-1/k\leq a$时，相应区间为空集。

反过来，任取$x\in(a,b)$，有$b-x>0$。选择正整数$k$满足$k\geq1/(b-x)$，则$1/k\leq b-x$，从而
\[a<x\leq b-1/k.\]
因此$x$属于右侧某一个区间，故$(a,b)$包含于该并集。两边互相包含，等式成立。
\end{proof}
\paragraph{答案}该可列并表示成立。

\subsection{作业第4题：无限重复试验的事件极限}\label{ex:w02-h4}
\paragraph{题目}
不断重复地做一个试验。设$0<p<1$，$A_n$表示前$n$次试验都不成功，且$P(A_n)=(1-p)^n$，$n\geq1$。
\begin{enumerate}[label=(\arabic*)]
\item 判断$A_n$是单增序列还是单减序列，并说明理由。
\item 说明$\lim_{n\to\infty}A_n$的含义。
\item 令$A$表示不断重复试验中至少成功一次，求$P(A)$。
\end{enumerate}
\basis{原笔记第1.2节事件的单调极限、概率的上连续性；第1.2节例“反复抽取扑克牌中的大王”。}
\paragraph{思路}先用事件的含义判断包含关系，再用概率上连续性。
\begin{proof}[解答]
\begin{enumerate}[label=(\arabic*)]
\item 前$n+1$次都不成功必然意味着前$n$次都不成功，所以$A_{n+1}\subseteq A_n$；这是单减事件列，不能只根据概率大小反推集合包含关系。
\item 单减列的极限为
\[\lim_{n\to\infty}A_n=\bigcap_{n=1}^{\infty}A_n,\]
表示每一次试验都不成功，即永远不成功。
\item 至少成功一次是永远不成功的补事件，因此
\begin{align*}
P(A)&=1-P\left(\bigcap_{n=1}^{\infty}A_n\right)\\
&=1-\lim_{n\to\infty}P(A_n)=1-\lim_{n\to\infty}(1-p)^n=1.
\end{align*}
使用上连续性是因为$A_n$单减且$P(A_1)\leq1<\infty$。概率为1不等于该事件在每个样本点都发生。
\end{enumerate}
\end{proof}
\paragraph{答案}单减；极限为永远不成功；至少成功一次的概率为1。
